%------------------------------------------------------------------------------%
\section{Demonstração por redução ao absurdo ou por contradição}

O método de redução ao absurdo, utiliza a negação da tese para chegar a uma
contradição da hipótese ou de alguma verdade matemática e é baseado na
\emph{lei do terceiro excluído} que diz que uma afirmação que não pode ser
falsa, deverá ser consequentemente verdadeira. Vamos ver na prática como essa
técnica funciona na prática provando a demonstração que segue:

%------------------------------------------------------------------------------%
\begin{proposicao}{Desigualdade das Médias}{prop:desigualdade_medias}
  Se $a$ e $b$ são números reais não-negativos, então
  \[
    \frac{a+b}{2} \geq \sqrt{ab}.
  \]
\end{proposicao}
%------------------------------------------------------------------------------%

%------------------------------------------------------------------------------%
\begin{proof}
  Inicialmente, vamos destacar a hipótese e a tese:
  \[
    \text{Hipótese: } a, b \in \R_{+}
  \]
  \[
    \text{Tese: } \frac{a+b}{2} \geq \sqrt{ab}
  \]
  Começamos a demonstração negando a tese até chegar em algum absurdo.
  Assim, vamos supor que \(\frac{a+b}{2} < \sqrt{ab}\).
  Desta maneira
  \begin{align*}
    \frac{a+b}{2} & < \sqrt{ab} \\
    a+b & < 2 \sqrt{ab} \\
    a-2\sqrt{ab}+b & <0 \\
    (\sqrt{a}-\sqrt{b})^2 & <0 \\
  \end{align*}
  Deste modo essa conclusão é um absurdo, já que nenhum número real ao quadrado
  pode ser negativo. Portanto,
  \[
    \frac{a+b}{2} < \sqrt{ab}
  \]
  é falsa e
  \[
    \frac{a+b}{2} \geq \sqrt{ab}
  \]
  é verdadeira.
\end{proof}
%------------------------------------------------------------------------------%

Vamos usar a redução ao absurdo para provar a seguinte proposição.

%------------------------------------------------------------------------------%
\begin{proposicao}{Irracionalidade de raiz de 2}{prop:raiz_2}
  O número $\sqrt{2}$ é irracional.
\end{proposicao}
%------------------------------------------------------------------------------%

%------------------------------------------------------------------------------%
\begin{proof}
  Vamos negar a tese, ou seja, vamos supor que $\sqrt{2}$ é racional. Isso
  significa dizer que podemos escrever $\sqrt{2}$ na forma de fração
  irredutível, isto é, $\sqrt{2}=\sfrac{p}{q}$, com $p$ e $q$ primos entre si
  (não possui divisores comuns além do 1). Dessa maneira
  \[
    \sqrt{2}=\frac{p}{q}
    \quad\Rightarrow\quad 2=\frac{p^2}{q^2}
    \quad\Rightarrow\quad p^2=2q^2
  \]
  onde a última igualdade nos diz que $p$ é um número par, em outras palavras
  $p=2k$, $k \in \Z$. Assim, retornando para a última igualdade
  \[
    p^2=2q^2
    \quad\Rightarrow\quad (2k)^2=2q^2
    \quad\Rightarrow\quad 4k^2=2q^2
    \quad\Rightarrow\quad q^2=2k^2
  \]
  Dessa maneira concluímos que $q$ também é um número par. Ora, sendo $p$ e $q$
  números pares, então $2$ é um divisor comum de $p$ e $q$, o que é um absurdo,
  visto que estamos supondo de $p$ e $q$ são primos entre si. Portanto, é um
  absurdo supor que $\sqrt{2}$ é racional e sendo assim, $\sqrt{2}$ é
  irracional.
\end{proof}
%------------------------------------------------------------------------------%

